% --- 12 Comparison with related work -----------------------------------------
% Provenance: figures/ch9-discussion/cross-architecture-table.tex (trimmed).
\begin{frame}{Σύγκριση με σχετικές εργασίες}
  \centering
  \renewcommand{\arraystretch}{1.2}
  {\small
  \begin{tabular}{@{}llll@{}}
    \toprule
    Εργασία & Πλατφόρμα & Προσέγγιση & Επιτάχυνση \\
    \midrule
    \en{Sch\"offel}~\cite{Schoeffel2022-HQC-IoT} & RV32 (SW) & καθαρό λογισμικό & — (baseline) \\
    \en{FAFFT}~\cite{OptHQC2025} & RV32 + \en{FPGA} & καθαρό λογισμικό & χαμηλή \\
    \en{pqm4}~\cite{Jang2025-CortexM4-HQC} & ARM Cortex-M4 & bit-slicing & μέτρια \\
    \midrule
    \HL{Παρούσα εργασία} & RV64GC + κρυπτο-\en{ISA} & \HL{εγγενείς εντολές} & \HL{έως \speedup{339}} \\
    \bottomrule
  \end{tabular}}
  \vspace{0.6em}
  \begin{itemize}
    \item Οι προηγούμενες RISC-V εργασίες είναι \con{καθαρά λογισμικές} ή απαιτούν custom \en{FPGA}
    \item \HL{Πρώτη εργασία} που αξιοποιεί \HL{επικυρωμένες κρυπτογραφικές επεκτάσεις} για το \HQC
  \end{itemize}
  \vspace{-0.2em}
  {\footnotesize Πλήρης πίνακας κύκλων/\instret στο παράρτημα.}
\end{frame}
